% =====================================================================
%  WSPÓLNE MAKRA PRACY
%  Oznaczenia zgodne z wspolne/slownik.md (sekcje 3.1 i 3.3).
%  Nowe makro dopisuje się tutaj, nie w pliku rozdziału.
% =====================================================================

% ---------------------------------------------------------------------
%  Znaczniki robocze
% ---------------------------------------------------------------------
% \TODO{tekst}  — czerwona wstawka w tekście + ostrzeżenie w logu
%                 („Package praca Warning: TODO: ...”).
% W wersji do oddania ustawić w main.tex \wersjaFinalnatrue:
% każdy pozostawiony \TODO zatrzyma kompilację błędem.
\newif\ifwersjaFinalna
\makeatletter
\newcommand{\TODO}[1]{%
  \ifwersjaFinalna
    \PackageError{praca}{Pozostawiony TODO: \detokenize{#1}}{Uzupełnij lub usuń znacznik.}%
  \else
    \PackageWarning{praca}{TODO: \detokenize{#1}}%
    \textcolor{red}{[TODO: #1]}%
  \fi}
\makeatother
% \DOSPRAWDZENIA{tekst} — fragment do weryfikacji (pomarańczowe tło znaczeniowe)
\newcommand{\DOSPRAWDZENIA}[1]{\textcolor{orange!80!black}{#1}}

% ---------------------------------------------------------------------
%  Odwołania (słownik, p. 5.5): małą literą, z tyldą
% ---------------------------------------------------------------------
\newcommand{\rysref}[1]{rys.~\ref{#1}}
\newcommand{\tabref}[1]{tab.~\ref{#1}}
\newcommand{\wzorref}[1]{wzór~\eqref{#1}}
\newcommand{\rozdzref}[1]{rozdz.~\ref{#1}}
\newcommand{\pref}[1]{p.~\ref{#1}}
\newcommand{\zalref}[1]{zał.~\ref{#1}}

% ---------------------------------------------------------------------
%  Terminy obcojęzyczne przy pierwszym użyciu: \ang{...} jest zajęte przez
%  siunitx (kąty), dlatego \angielski.
%  „spektroskopia EDS (\angielski{Energy Dispersive X-ray Spectroscopy})”
% ---------------------------------------------------------------------
\newcommand{\angielski}[1]{ang.~\textit{#1}}

% ---------------------------------------------------------------------
%  Rysunki i podpisy
% ---------------------------------------------------------------------
% \rysunek[szerokość]{plik}{podpis}{etykieta}
%   \rysunek[0.8\linewidth]{rozdz4/stanowisko.jpg}{Stanowisko pomiarowe
%     \opracowanieWlasne}{fig:stanowisko}
\newcommand{\rysunek}[4][\linewidth]{%
  \begin{figure}[htbp]
    \centering
    \includegraphics[width=#1]{#2}
    \caption{#3}
    \label{#4}
  \end{figure}}
% Dopiski o źródle rysunku (słownik, p. 5.10)
\newcommand{\opracowanieWlasne}{(opracowanie własne)}
\newcommand{\naPodstawie}[1]{(na podstawie~\cite{#1})}
\newcommand{\zrodlo}[1]{(źródło:~\cite{#1})}

% Tabela: podpis nad tabelą (ustawione w preambule), wzór:
% \begin{table}[htbp]\centering
%   \caption{Podpis}\label{tab:xxx}
%   \begin{tabular}{@{}lr@{}}\toprule A & B\\\midrule ... \\\bottomrule\end{tabular}
% \end{table}
% Nagłówek kolumny z jednostką: „$U_\mathrm{B}$ [\si{\volt}]” albo \tabjedn
\newcommand{\tabjedn}[2]{#1~[\si{#2}]}

% \inputBezPreambuly{plik} — wstawia fragment przygotowany jako samodzielny
% dokument (\documentclass, \usepackage, \sisetup, \begin{document} są
% pomijane; brak \end{document} nie przeszkadza). Przykład:
%   \resizebox{\textwidth}{!}{\inputBezPreambuly{../uklad-pomiarowy/uklad_pomiarowy_body}}
\makeatletter
\let\praca@begin\begin
\newcommand{\inputBezPreambuly}[1]{\begingroup
  \renewcommand{\documentclass}[2][]{}%
  \renewcommand{\usepackage}[2][]{}%
  \renewcommand{\sisetup}[1]{}%
  \def\begin##1{\ifstrequal{##1}{document}{}{\praca@begin{##1}}}%
  \input{#1}\endgroup}
\makeatother

% Domyślne katalogi z rysunkami i danymi dla fragmentów, które używają
% \figdir / \datadir (np. emisja-magnes-pomiary/emisja_tresc.tex).
% Wrapper rozdziału ustawia je przez \renewcommand przed \input.
\newcommand{\figdir}{rysunki}
\newcommand{\datadir}{dane}

% ---------------------------------------------------------------------
%  Oznaczenia wielkości (słownik, p. 3.1) — można też pisać wprost
% ---------------------------------------------------------------------
\newcommand{\UB}{U_\mathrm{B}}          % napięcie ekstrakcyjne (emiter–bramka)
\newcommand{\UT}{U_\mathrm{T}}          % napięcie przyspieszające (emiter–target)
\newcommand{\Ie}{I_\mathrm{e}}          % prąd emisji
\newcommand{\IB}{I_\mathrm{B}}          % prąd bramki
\newcommand{\IT}{I_\mathrm{T}}          % prąd targetu
\newcommand{\RT}{R_\mathrm{T}}          % rezystor w obwodzie targetu
\newcommand{\RB}{R_\mathrm{B}}          % rezystor w obwodzie bramki
\newcommand{\Ezero}{E_0}                % energia elektronów wiązki
\newcommand{\Ec}{E_\mathrm{c}}          % energia krytyczna
\newcommand{\me}{m_\mathrm{e}}          % masa elektronu
\newcommand{\kB}{k_\mathrm{B}}          % stała Boltzmanna
\newcommand{\FL}{\vec F_\mathrm{L}}     % siła Lorentza
\newcommand{\RKO}{R_\mathrm{KO}}        % zasięg Kanayi–Okayamy
\newcommand{\dd}{\mathop{}\!\mathrm{d}} % różniczka: \dd x

% Linie widmowe (notacja Siegbahna): „Si~\Ka”, „Fe~\Kb”
\newcommand{\Ka}{K\ensuremath{\alpha}}
\newcommand{\Kb}{K\ensuremath{\beta}}
\newcommand{\La}{L\ensuremath{\alpha}}

% Jednostka umowna (jasność ekranu L)
\DeclareSIUnit{\ju}{j.u.}

% Plik z dodatkowymi makrami wątków (opcjonalny)
\InputIfFileExists{szablon/makra_dodatkowe.tex}{}{}
